\documentclass[10.5pt]{article}
% Shared style for the Airspace Monitor reports
\usepackage[a4paper,margin=1.9cm,top=2.25cm,bottom=2.05cm,headheight=18pt]{geometry}
\usepackage{fontspec}
\usepackage{xcolor}
\usepackage{graphicx}
\usepackage{tikz}
\usetikzlibrary{positioning,calc,shapes.geometric,backgrounds,arrows.meta}
\usepackage[most]{tcolorbox}
\usepackage{booktabs}
\usepackage{amssymb}
\usepackage{tabularx}
\usepackage{longtable}
\usepackage{array}
\usepackage{colortbl}
\usepackage{ragged2e}
\usepackage{enumitem}
\usepackage{fancyhdr}
\usepackage{fontawesome5}
\usepackage{hyperref}
\usepackage{microtype}
\usepackage{lastpage}

\definecolor{ink}{HTML}{0B1020}
\definecolor{panel}{HTML}{111A2E}
\definecolor{border}{HTML}{22304F}
\definecolor{textlight}{HTML}{E7ECF6}
\definecolor{muted}{HTML}{9AA8C2}
\definecolor{accent}{HTML}{4CC3FF}
\definecolor{violet}{HTML}{B392FF}
\definecolor{okgreen}{HTML}{2FAE78}
\definecolor{okgreenlight}{HTML}{E6F7F0}
\definecolor{warn}{HTML}{C77A00}
\definecolor{warnlight}{HTML}{FFF6E5}
\definecolor{alert}{HTML}{C94646}
\definecolor{alertlight}{HTML}{FDEEEE}
\definecolor{paper}{HTML}{F7F9FD}
\definecolor{papertext}{HTML}{1B2438}
\definecolor{papermuted}{HTML}{5B6B85}
\definecolor{lightborder}{HTML}{DCE4F2}
\definecolor{softblue}{HTML}{EAF6FF}
\definecolor{softviolet}{HTML}{F1ECFF}
\definecolor{softgrey}{HTML}{F0F3F9}

\setmainfont{Lato}
\setsansfont{Lato}
\setmonofont{Consolas}[Scale=0.92]
\renewcommand{\familydefault}{\sfdefault}
\setlength{\parindent}{0pt}
\setlength{\parskip}{5pt}
\RaggedRight
\renewcommand{\arraystretch}{1.32}
\setlength{\tabcolsep}{5pt}
\setlength{\emergencystretch}{1em}
\hypersetup{hidelinks}

\newcommand{\sechead}[1]{%
  \vspace{7pt}%
  {\color{papertext}\fontsize{17}{20}\selectfont\bfseries #1}\par
  \vspace{-2pt}{\color{accent}\rule{\linewidth}{1.3pt}}\par
  \vspace{5pt}%
}
\newcommand{\subhead}[1]{%
  \vspace{5pt}{\color{papertext}\fontsize{11.5}{14}\selectfont\bfseries #1}\par
  \vspace{1pt}%
}
\newcommand{\code}[1]{\texttt{\detokenize{#1}}}
\newcommand{\docsrc}[1]{{\scriptsize\color{papermuted}\texttt{\detokenize{#1}}}}

\newcommand{\pill}[2]{%
  \begin{tikzpicture}[baseline=(x.base)]%
    \node[fill=#1,text=white,rounded corners=2.4pt,inner sep=2.6pt,font=\scriptsize\bfseries] (x) {#2};%
  \end{tikzpicture}%
}
\newcommand{\dotok}{\textcolor{okgreen}{\faCircle}}
\newcommand{\dotwarn}{\textcolor{warn}{\faCircle}}
\newcommand{\dotbad}{\textcolor{alert}{\faCircle}}

\newtcolorbox{card}[1][]{enhanced,breakable,colback=white,colframe=lightborder,boxrule=0.6pt,arc=2.2mm,left=4mm,right=4mm,top=3mm,bottom=3mm,#1}
\newtcolorbox{accentcard}[1][]{enhanced,breakable,colback=softblue,colframe=accent,boxrule=0.8pt,arc=2.2mm,left=4mm,right=4mm,top=3mm,bottom=3mm,#1}
\newtcolorbox{warncard}[1][]{enhanced,breakable,colback=warnlight,colframe=warn,boxrule=0.8pt,arc=2.2mm,left=4mm,right=4mm,top=3mm,bottom=3mm,#1}
\newtcolorbox{badcard}[1][]{enhanced,breakable,colback=alertlight,colframe=alert,boxrule=0.8pt,arc=2.2mm,left=4mm,right=4mm,top=3mm,bottom=3mm,#1}
\newtcolorbox{violetcard}[1][]{enhanced,breakable,colback=softviolet,colframe=violet,boxrule=0.8pt,arc=2.2mm,left=4mm,right=4mm,top=3mm,bottom=3mm,#1}
\newtcolorbox{softcard}[1][]{enhanced,breakable,colback=softgrey,colframe=lightborder,boxrule=0.6pt,arc=2.2mm,left=4mm,right=4mm,top=3mm,bottom=3mm,#1}

\pagestyle{fancy}
\fancyhf{}
\fancyhead[L]{\scriptsize\color{papermuted}Airspace Monitor · Documentation Review}
\fancyhead[R]{\scriptsize\color{papermuted}30 October 2026}
\fancyfoot[L]{\scriptsize\color{papermuted}Level 1 · pre-implementation}
\fancyfoot[R]{\scriptsize\color{papermuted}\thepage\ / \pageref{LastPage}}
\renewcommand{\headrulewidth}{0.4pt}
\renewcommand{\footrulewidth}{0.4pt}
\renewcommand{\headrule}{\hbox to\headwidth{\color{lightborder}\leaders\hrule height \headrulewidth\hfill}}
\renewcommand{\footrule}{\hbox to\headwidth{\color{lightborder}\leaders\hrule height \footrulewidth\hfill}}


% Summary-specific header/footer
\pagestyle{fancy}
\fancyhf{}
\fancyhead[L]{\footnotesize\color{papermuted}Airspace Monitor · Project summary}
\fancyhead[R]{\footnotesize\color{papermuted}2026}
\fancyfoot[L]{\footnotesize\color{papermuted}Data: The OpenSky Network · non-commercial use}
\fancyfoot[R]{\footnotesize\color{papermuted}\thepage\ / 2}
\renewcommand{\headrulewidth}{0.4pt}
\renewcommand{\footrulewidth}{0.4pt}

\newtcolorbox{darkcard}[1][]{enhanced,colback=panel,colframe=border,coltext=textlight,boxrule=0.6pt,arc=2.2mm,left=3.3mm,right=3.3mm,top=2.6mm,bottom=2.6mm,#1}

\begin{document}
\hypersetup{pdftitle={Airspace Monitor - Project Summary},pdfauthor={Alek Hanusic and Daniel Csorba},pdfsubject={A plain-language summary of the Airspace Monitor project}}

% ============================================================
% PAGE 1 — WHAT IT IS
% ============================================================
\pagecolor{ink}
\color{textlight}
\thispagestyle{empty}

\begin{tikzpicture}[remember picture,overlay]
  \foreach \x in {0,...,14} {
    \draw[white,opacity=0.035,line width=0.4pt]
      ($(current page.south west)+(\x*1.55cm,0)$) -- ++(0,29.7cm);
  }
  \foreach \y in {0,...,19} {
    \draw[white,opacity=0.035,line width=0.4pt]
      ($(current page.south west)+(0,\y*1.55cm)$) -- ++(21cm,0);
  }
  \draw[accent,opacity=0.30,line width=0.8pt]
    ($(current page.north east)+(-4.6cm,-3.2cm)$) arc[start angle=205,end angle=515,radius=2.0cm];
  \draw[accent,opacity=0.14,line width=0.8pt]
    ($(current page.north east)+(-4.6cm,-3.2cm)$) arc[start angle=205,end angle=515,radius=2.8cm];
  \node[circle,draw=accent,fill=accent!12,minimum size=1.0cm,line width=0.8pt]
    at ($(current page.north east)+(-4.6cm,-3.2cm)$) {\color{accent}\large\faPlane};
\end{tikzpicture}

\vspace*{1.35cm}
{\color{accent}\fontsize{9}{11}\selectfont\bfseries MSc CLOUD \& AI PROJECT \quad·\quad 2026}

\vspace{0.35cm}
{\fontsize{42}{46}\selectfont\bfseries Airspace Monitor}

\vspace{0.10cm}
{\color{textlight}\fontsize{20}{25}\selectfont See the sky. Understand the traffic. Ask the airspace.}

\vspace{0.55cm}
{\color{muted}\fontsize{11.2}{15.5}\selectfont
Airspace Monitor is a live map of the aircraft visible to the OpenSky volunteer network. It turns free, crowdsourced flight data into a real-time picture of the sky — and adds predictions and an AI assistant on top.\par}

\vspace{0.55cm}
\begin{minipage}[t]{0.485\linewidth}
\begin{darkcard}
{\color{accent}\Large\faPlane}\quad{\bfseries Live map}

\vspace{2pt}
{\small 10,000+ aircraft move smoothly between updates, coloured by altitude. Click one to see its flight, route and predicted arrival.}
\end{darkcard}
\end{minipage}\hfill
\begin{minipage}[t]{0.485\linewidth}
\begin{darkcard}
{\color{violet}\Large\faLayerGroup}\quad{\bfseries Unusual traffic}

\vspace{2pt}
{\small Hexagonal map cells show where traffic is busier or quieter than usual for that hour of the week — with an honest caveat when coverage may be the cause.}
\end{darkcard}
\end{minipage}

\vspace{4pt}
\begin{minipage}[t]{0.485\linewidth}
\begin{darkcard}
{\color{accent}\Large\faClock}\quad{\bfseries Arrival board}

\vspace{2pt}
{\small Predicted landing times, shown as a range, for flights inbound to Vienna, Bratislava and Budapest.}
\end{darkcard}
\end{minipage}\hfill
\begin{minipage}[t]{0.485\linewidth}
\begin{darkcard}
{\color{violet}\Large\faComments}\quad{\bfseries Ask the sky}

\vspace{2pt}
{\small A plain-English chat that queries the database and shows the SQL behind each answer, plus a short AI briefing every 30 minutes.}
\end{darkcard}
\end{minipage}

\vspace{8pt}
\begin{center}
\begin{tabular}{@{}p{3.4cm} p{3.4cm} p{3.4cm} p{3.4cm}@{}}
{\color{accent}\fontsize{20}{23}\selectfont\bfseries 10k+} & {\color{accent}\fontsize{20}{23}\selectfont\bfseries 45--90 s} & {\color{accent}\fontsize{20}{23}\selectfont\bfseries 3} & {\color{accent}\fontsize{20}{23}\selectfont\bfseries 30 min} \\
{\color{muted}\scriptsize aircraft in one snapshot} & {\color{muted}\scriptsize refresh between updates} & {\color{muted}\scriptsize target airports} & {\color{muted}\scriptsize AI briefing cadence} \\
\end{tabular}
\end{center}

\vspace{10pt}
{\color{accent}\fontsize{8.5}{10}\selectfont\bfseries WHY IT IS INTERESTING}

\vspace{3pt}
\begin{minipage}[t]{0.32\linewidth}
\begin{darkcard}[colframe=accent]
{\bfseries Real data engineering}

\vspace{2pt}
{\footnotesize Two collectors spool every response before upload. Ingest is idempotent, raw data is archived to R2, and the restore path is documented.}
\end{darkcard}
\end{minipage}\hfill
\begin{minipage}[t]{0.32\linewidth}
\begin{darkcard}[colframe=accent]
{\bfseries Honest ML}

\vspace{2pt}
{\footnotesize Arrival times are judged against two baselines and returned as a range. Unusual traffic is evaluated with synthetic anomalies because there is no ground truth.}
\end{darkcard}
\end{minipage}\hfill
\begin{minipage}[t]{0.32\linewidth}
\begin{darkcard}[colframe=accent]
{\bfseries Guarded AI}

\vspace{2pt}
{\footnotesize The chat runs read-only SQL against curated views, shows the query, and is tested against a 30-question golden set with an 80\% target.}
\end{darkcard}
\end{minipage}

\vspace{4pt}
{\color{muted}\fontsize{9}{12}\selectfont\itshape
A map you can look at, a platform you can inspect, and predictions you can question.}

\vfill
{\color{muted}\fontsize{8.5}{11}\selectfont
Free, non-commercial data. The OpenSky attribution is always visible in the app.\\
Built by Alek Hanušić \& Dániel Csorba.\par
\vspace{2pt}
{\color{accent}\rule{1.1cm}{0.9pt}}\quad
\textbf{A live map with a data platform, predictions and a guarded AI assistant underneath.}
}

\newpage
\pagecolor{paper}
\color{papertext}

% ============================================================
% PAGE 2 — HOW IT WORKS
% ============================================================
\sechead{How it works}

{\small
Four steps, all built as one system:
}

\vspace{2pt}
\begin{minipage}[t]{0.235\linewidth}
\begin{card}
{\color{accent}\bfseries 1 · Collect}

\vspace{2pt}
{\footnotesize Two home PCs ask OpenSky for every visible aircraft every 90 seconds, taking turns. Each response is saved locally first, so a network or server outage only delays data; it does not lose it.}
\end{card}
\end{minipage}\hfill
\begin{minipage}[t]{0.235\linewidth}
\begin{card}
{\color{accent}\bfseries 2 · Store}

\vspace{2pt}
{\footnotesize The cloud server stores positions in PostgreSQL with TimescaleDB and PostGIS. Older raw data is archived to Cloudflare R2 and read with DuckDB when needed.}
\end{card}
\end{minipage}\hfill
\begin{minipage}[t]{0.235\linewidth}
\begin{card}
{\color{accent}\bfseries 3 · Understand}

\vspace{2pt}
{\footnotesize Scheduled jobs build flights, detect squawks and go-arounds, score unusual traffic, look up likely destinations and predict arrival times with a LightGBM model.}
\end{card}
\end{minipage}\hfill
\begin{minipage}[t]{0.235\linewidth}
\begin{card}
{\color{accent}\bfseries 4 · Show}

\vspace{2pt}
{\footnotesize A React map streams new snapshots over SSE. The AI briefing and chat read curated database views; the chat can only run read-only SQL and always shows the query.}
\end{card}
\end{minipage}

\vspace{5pt}
\sechead{What is inside}

\begin{minipage}[t]{0.32\linewidth}
\begin{softcard}
{\bfseries Data platform}

\vspace{2pt}
{\footnotesize Python, FastAPI, PostgreSQL 16 + TimescaleDB + PostGIS, Dagster, Grafana, Cloudflare R2. Idempotent ingest, spool-first collectors, data-quality checks and a documented restore path.}
\end{softcard}
\end{minipage}\hfill
\begin{minipage}[t]{0.32\linewidth}
\begin{softcard}
{\bfseries Machine learning}

\vspace{2pt}
{\footnotesize A statistical unusual-traffic baseline, a LightGBM arrival-time model with p10/p50/p90 outputs, and a deterministic route lookup. Every model is compared with a simple baseline and has a public model card.}
\end{softcard}
\end{minipage}\hfill
\begin{minipage}[t]{0.32\linewidth}
\begin{softcard}
{\bfseries AI layer}

\vspace{2pt}
{\footnotesize Claude Haiku 5.5 writes the briefing from a facts JSON and answers chat questions through a guarded text-to-SQL path: read-only role, 5 s timeout, allow-listed views and functions, 3 tool calls, 200 rows, SQL shown.}
\end{softcard}
\end{minipage}

\vspace{4pt}
\begin{minipage}[t]{0.55\linewidth}
\begin{warncard}
{\bfseries What it cannot do (and says so)}
\begin{itemize}[leftmargin=1.0em,itemsep=1.5pt,topsep=2pt]
  \item It only sees aircraft heard by OpenSky's volunteer receivers. A receiver outage can look like a traffic drop; low-traffic cells carry a coverage caveat.
  \item There are no airline schedules or air-traffic-control data. Arrival times are ranges, and far-out estimates are wider.
  \item Destinations are estimates, not broadcast data. Diversions can appear late.
  \item It is not a commercial product, a mobile app or a user-account service.
\end{itemize}
\end{warncard}
\end{minipage}\hfill
\begin{minipage}[t]{0.42\linewidth}
\begin{accentcard}
{\bfseries Status}

\vspace{2pt}
{\small Week 0 of 8. The plan and technical documents are finished; collection starts first so the models have history.}

\vspace{2pt}
{\small \textbf{Core:} live map, unusual-traffic heatmap, event feed, arrival board, destination estimate, AI briefing and chat, public dashboards.}

\vspace{2pt}
{\small \textbf{Stretch, only if time:} a smarter destination model, a smarter anomaly model, holding-pattern events, self-hosted map tiles.}
\end{accentcard}
\end{minipage}

\vspace{4pt}
{\footnotesize\color{papermuted}
\textbf{Data and credits:} The OpenSky Network (live states and flight lists), Zenodo OpenSky COVID-19 flight dataset, OurAirports, OpenFlights, Natural Earth, aviationweather.gov and the Iowa Environmental Mesonet. Basemap © CARTO © OpenStreetMap contributors. OpenSky data is used for non-commercial research and education; attribution is always visible.
}

\end{document}
